\documentclass[10pt,twocolumn]{article}
\usepackage[a4paper,margin=15mm]{geometry}
\usepackage{amsmath,amssymb}
\usepackage{mathrsfs}
\usepackage{amsthm}
\theoremstyle{definition}
\newtheorem{remark}{Remark}
\usepackage{graphicx}
\usepackage{booktabs,array,calc}
\usepackage[font=small,labelfont=bf]{caption}
\usepackage[colorlinks=true]{hyperref}
\usepackage{etoolbox}
\providecommand{\tightlist}{\setlength{\itemsep}{0pt}\setlength{\parskip}{0pt}}
\setcounter{secnumdepth}{-1}
\emergencystretch=3em
\allowdisplaybreaks

\begin{document}
\title{From Obstruction Certificates to Monitoring Standards:\\ An Applications and Limitations Statement\\ for the Obstruction Calculus}
\author{Amin Abaee\\[0.35em]
{\small Independent Researcher}\\[0.55em]
{\small\href{https://orcid.org/0000-0002-0019-1842}{ORCID: 0000-0002-0019-1842}}\\[0.3em]
{\small\href{mailto:amin\_abaee@ut.ac.ir}{amin\_abaee@ut.ac.ir}}}
\date{September 24, 2026}
\maketitle

\begin{abstract}
This statement consolidates, in one place and in plain scope, what the
obstruction calculus and its follow-up programme are for, whom they serve
today, and what they do not yet do. It is a satellite document: it creates
no new theorems and modifies no manuscript of the programme; every claim
below is a pointer into a cited source with a machine-verifiable record.
Four design rules are extracted in usable form --- review-interval sizing,
the fibre rule for indicator aggregation, bias correction, and
structural-cause forensics --- each with its scope, its source theorem,
and its verified instance. The operational-versus-architectural
distinction is stated explicitly, a stakeholder scope table is given, and
the limitations are recorded with the same precision as the
capabilities: model-availability, the dimensionality wall of exact
recursions, the deferred numerical campaign, the absence of calibrated
case studies, and the verification-grade status of the software. The
programme's delivered lineages are indexed against the utility
prerequisites, and the genuinely open items are listed.
\end{abstract}

\subsection{1. Purpose and status}\label{purpose}

The obstruction calculus (Abaee, 2026a) proves, from a model of a
controlled system and its observation structure, that certain information
deficits --- coarse aggregation, delayed review, biased reading,
unavailable certification --- make robust viability impossible even when
every physical state is viable in isolation. The follow-up programme has
carried the calculus to a continuous-to-finite bridge with solver-level
certificates (Abaee, 2026b), a selector unification (Abaee, 2026c), an
exact stochastic belief-state value theory (Abaee, 2026d), a five-layer
successor architecture (Abaee, 2026e), and a generalization catalogue of
ten further lineages (Abaee, 2026f). This statement is a satellite of the
whole set: it translates the results into the form in which they can be
used --- and states, with the same care, the boundary of that use. Nothing
here creates theorem status; every pointer resolves to a source with an
exact-arithmetic verification record, and the pointers themselves are
machine-checked (Section 7).

\subsection{2. Four design rules}\label{rules}

\paragraph{R1 (review-interval sizing).}
Within a declared hold-until-review class --- the institutional reading of
a sampled-review regulator --- a review period exceeding the worst-case
exit time renders in-window enforcement futile: the timing certificate
(Theorem 3 of Abaee, 2026a) certifies nonviability exactly when
\(T_{\mathrm{obs}}\) exceeds the exit-time value. The scope qualifier is
load-bearing: the statement is about the declared class, and the
stochastic companion quantifies the contrast --- on the audited
delayed-regime grid the declared-class value and the unrestricted
sequential value coincide at the one-step horizon and differ exactly on
the timing cells for later horizons, where an open-loop time-varying
sequence keeps both regime branches within \([z_{0}-1, z_{0}+1]\)
(Abaee, 2026d, Theorem~2(e) and Remark on class scope; verified by the
oscillation identity, check S11, and the class table, check S8). On the
audited grid the exit-time value is \(\sigma^{*}(z_{0}) = z_{0} - 1\)
exactly (48 cells; Abaee, 2026d, checks S4/S15). As a design rule: size
the review period against the worst-case class-scoped exit time, and
state the class when invoking it.

\paragraph{R2 (the fibre rule for indicator aggregation).}
An aggregated indicator can support an exact safe/unsafe verdict on a
floor only if its fibres do not merge states requiring incompatible
actions: no threshold on a composite index can guarantee a specific
coordinate floor in general (Proposition 7 and Corollary 1 of Abaee,
2026a, which apply verbatim to composite indices). The rule has an exact
one-dimensional design form --- every observation fibre must carry an
instrument jointly safe for all its states --- and its cost reading is
computed on the audited target: the coarsest admissible monitoring costs
two cells, and the pairwise version of the rule is necessary but not
sufficient off convexity, with an exact three-codex counterinstance
(Abaee, 2026f, lineage \emph{From Obstruction to Design}, checks E1--E3). As a
design rule: coordinate-wise floor monitoring is the sound default; an
aggregate is admissible only with certified compatible fibres.

\paragraph{R3 (bias correction).}
A known reading bias empties the kernel of an uncorrected
certainty-equivalence controller even when the perfect-information kernel
is nonempty (Remark 4 of Abaee, 2026a). In the audited trap the
uncorrected reading drives every class law to drift at least
\(\tfrac{21}{100}\) --- exit --- on the whole audited band, while the
corrected law is viable (Abaee, 2026f, lineage \emph{Beyond Constant Drift};
Abaee, 2026d, check S1's two-floor instance for the corresponding chance
bound). As a design rule: a biased indicator is admissible only under a
correction certified against the model, not merely acknowledged.

\paragraph{R4 (structural-cause forensics).}
The certificates separate two post-mortem narratives that are otherwise
conflated: managerial failure and information-structure impossibility.
The common-action obstruction certifies that the reading merged states
requiring opposite interventions --- collapse guaranteed before any
decision --- with exact instances throughout the programme (the audited
belief \(\{(1,2),(2,1)\}\): nonviable under the aggregate reading,
viable under either coordinate reading; Abaee, 2026d, check S5; the
unanimity deadlock of decentralized observation, sixteen coordinated-
viable pairs lost under split rights, Abaee, 2026f, lineage \emph{Decentralized Observation}). As a design
rule: a collapse audit should test the information structure before the
personnel.

\begin{remark}[what the rules are]
Each rule is a restatement of a theorem at design precision, with its
scope and its verified instance attached. None is a heuristic; none
extends its source. Where a rule is class-scoped (R1), the scope is part
of the rule.
\end{remark}

\subsection{3. Operational versus architectural use}\label{levels}

The calculus is not an operational control tool: it does not steer a
fishery, a reservoir, or a vehicle, and no part of the programme claims
otherwise. Its use is architectural and forensic: (i) \emph{ex-ante
design} --- proving a proposed monitoring arrangement (aggregation,
review period, bias treatment, sensor suite) mathematically unable to
support safety, before it is built; (ii) \emph{standard-setting} ---
converting qualitative objections to composite-index compliance into
impossibility theorems with the precision required in regulatory texts;
(iii) \emph{forensics} --- separating structural from managerial cause
after failure; and (iv) \emph{sensor-sufficiency certification} in
safety-critical automation --- certifying a degraded sensing configuration
inadequate against worst-case disturbance, the direction opposite to
conventional safety verification, with the programme's policy-class and
hybrid-mode lineages as the natural instruments (Abaee, 2026f, lineages \emph{The Policy-Class Lattice} and \emph{Discrete Regimes, Carefully}).

\subsection{4. Stakeholder scope}\label{stakeholders}

\begin{table*}[t]
\centering
\footnotesize
\begin{tabular}{@{}l p{0.30\textwidth}@{}}
\toprule
Stakeholder & Reach today \\
\midrule
Field operator & none direct; rules R1--R3 enter as institutional \\
Standards body & R1--R2 as auditable design constraints \\
Post-mortem auditor & R4 as a structural-cause test \\
Safety engineer (CPS) & sensor-sufficiency certification (outlook) \\
Viability/control researchers & the full calculus and programme \\
\bottomrule
\end{tabular}
\caption{Who can use the results today, and at which level.}
\end{table*}

\subsection{5. Limitations, stated precisely}\label{limits}

\paragraph{L1 (model availability).}
The calculus consumes a model: dynamics, disturbance bounds, floor
boundaries, instrument sets. Ecological and institutional systems rarely
arrive with calibrated equations and tight bounds. The answer is not to
shrink the requirement but to name it: \emph{bounding is the regulatory
act}. Declaring the disturbance class and the floor margins is precisely
the deliberation the calculus formalizes, and the buffered-viability
lineage prices margins into the certificates (Abaee, 2026f, lineage \emph{Buffered Viability}: the
obstruction margin, e.g.\ \(\mu = 2\) on the audited beliefs, erosion
\(28 \to 3 \to 0\)). Where no defensible bounds exist, the calculus is
silent --- by construction, not by oversight.

\paragraph{L2 (dimensionality).}
Exact recursions are exponential in the worst case
(\(|\Gamma_{k+1}| \le |A|\,|\Gamma_{k}|^{|Y|}\); Abaee, 2026d,
Proposition~1); the audited instances are one- and two-dimensional, and
ten-to-thirty-variable ecosystem models are not auditable by the exact
route. The continuous-to-finite bridge carries LP-level certificates for
polyhedral classes (Abaee, 2026b, solver-certified campaign), the
timing-instantiation theorem reduces the timing check to finite linear
programs over polytope-declared classes (Abaee, 2026a, Theorem~4), and
the finite-horizon completeness lineage gives an exact game algorithm on
finite systems (Abaee, 2026f, lineage \emph{Complete on the Finite Horizon}). A medium-dimensional numerical campaign
remains deferred in the calculus manuscript and is the main computational
gap.

\paragraph{L3 (completeness scope).}
Completeness of the obstruction calculus holds on finite horizons for
finite systems (Abaee, 2026f, lineage \emph{Complete on the Finite Horizon}: soundness, completeness, and
counterstrategy trees; the certificate gap is empty there); for
continuous blind-window control classes the check remains a template
(Open Problem~1 of Abaee, 2026a), and the residual dynamic gap of its
\(\S\)6.5 stands. Statements beyond these scopes are not claimed.

\paragraph{L4 (calibration and case studies).}
No result in the programme is calibrated against real schedules, stocks,
or sensor data. The two-patch, hidden-regime, and trap instances are
deliberately small and exact; the real-schedule illustration in the
calculus manuscript was marked optional and not taken. Applied
calibration is an open, owner-side item.

\paragraph{L5 (software status).}
The programme's software is verification-grade: exact rational scripts
regenerating every claimed number (sixty check families in the catalogue alone; fifteen in the stochastic lineage), plus a solver-certified LP
campaign. There is no user-facing decision-support tool; the distance
from a verification script to a tool is an engineering programme of its
own and is recorded as such.

\paragraph{L6 (formalizing the intuitive).}
That aggregation hides trade-offs, that late review arrives too late,
that uncorrected bias misleads --- all intuitive. The contribution is that
these are theorems: in standard-setting and forensic contexts, intuition
is not admissible evidence, and an impossibility proof with a verified
certificate is. The calculus' value is not that its conclusions surprise
experts; it is that its conclusions bind.

\subsection{6. The delivered programme against the utility prerequisites}\label{delivered}

The follow-up programme was drafted against exactly the prerequisites an
operational reading would set. Index: continuous-to-finite checkable
certificates --- delivered (Abaee, 2026b, exact primal/dual LP witnesses);
belief-state value theory --- delivered (Abaee, 2026d, fifteen exact
check families); monitoring design rules --- delivered (Abaee, 2026f, lineage \emph{From Obstruction to Design}:
fibre conditions, coarsest monitoring, delay/bias/aggregation rules);
diagnostics for restriction attribution --- delivered (Abaee, 2026f, lineage \emph{The Policy-Class Lattice}: computed distances to viability);
exit-time theory beyond constant drift --- delivered (Abaee, 2026f, lineage \emph{Beyond Constant Drift});
uncertainty, regimes, margins, and institutional splits --- delivered
(Abaee, 2026f, lineages \emph{Hidden Parameters and Learning Deadlines} through \emph{Decentralized Observation}); architecture consolidating the five layers ---
delivered (Abaee, 2026e). Remaining: user-facing tooling (L5), applied
calibration (L4), the deferred numerical campaign (L2), and the open
completeness problems (L3).

\subsection{7. Machine-checked pointers}\label{pointers}

\texttt{paper2\_applications\_note\_v1\_pointers.py} probes the cited
artifacts (the shipped PDFs of the programme) for each object this
statement quotes --- the theorems, the verified instances, and the
lineage titles --- and reports the pointer record. It requires a PDF text
extractor (\texttt{pymupdf}, stated dependency); all other checks are
exact string identities on normalized text. A pointer that fails breaks
the statement; the record ships with it.

\subsection{References}\label{references}
{\footnotesize

Abaee, A.: An obstruction calculus for viability under incomplete
observation. Preprint.
\url{https://github.com/MIKEAA2020/general-sustainability} (2026a)

Abaee, A.: Finite nonviability certificates under delayed information.
Preprint. \url{https://github.com/MIKEAA2020/general-sustainability} (2026b)

Abaee, A.: The viable selector: a unifying framework for obstruction
certificates under incomplete observation. Preprint.
\url{https://github.com/MIKEAA2020/general-sustainability} (2026c)

Abaee, A.: The stochastic selector: exact rational belief-state safety
values under partial observation, 2nd edn. Preprint.
\url{https://github.com/MIKEAA2020/general-sustainability} (2026d)

Abaee, A.: Five layers of obstruction: a successor architecture for
viability under incomplete observation. Preprint.
\url{https://github.com/MIKEAA2020/general-sustainability} (2026e)

Abaee, A.: Generalization catalogue of the obstruction programme, ten
lineages: vector floors; finite-horizon completeness; certificate
duality; the policy-class lattice; monitoring design; comparison-function
exit; hidden parameters and learning deadlines; hybrid modes; buffered
viability; decentralized observation. Preprint.
\url{https://github.com/MIKEAA2020/general-sustainability} (2026f)

}

\section*{Declarations}

\textbf{Funding.} None.\quad \textbf{Competing interests.} None.\quad
\textbf{Data availability.} No external data were used.\quad
\textbf{Code availability.} The pointer-check script
\texttt{paper2\_applications\_note\_v1\_pointers.py} is publicly available
at \url{https://github.com/MIKEAA2020/general-sustainability} (folder
\texttt{arena agent 1/paper rewrites/latex} in the repository); it
verifies every pointer of Section 7 against the shipped
artifacts.\quad
\textbf{AI declaration.} GLM (Z.ai), Qwen (Alibaba Cloud) and DeepSeek AI
assisted with drafting and iterative review within the programme's
verification discipline.

\end{document}
